これにより、簡単のためチルダを省略すると、次を得る。
\OriginalEquation{179}
\(\pi\) に変換則
\OriginalEquation{180}
を与えれば、上の作用が、無限小変換に限らず、あらゆる次数で微分同相変換に対して不変であることを確かめられる。この定義では、時間の微分同相変換のもとで、\(\pi\) はスカラー場としての変換に加えてシフトを受ける。もともと微分同相不変な項には \(\pi\) が現れないことに注意しよう。

この手続きをユニタリーゲージの作用\eqref{eq:163}に適用すると、
\OriginalEquation{181}
を得る。ここではひとまず、簡単のため外部曲率を含む項を無視した。

この作用はかなり複雑であり、この段階では、ユニタリーゲージのラグランジアンからGoldstoneボソン \(\pi\) を再導入する利点は明らかではない。ゲージ理論の場合と同様、十分短い距離では計量の揺らぎを無視してGoldstoneボソンの物理を調べられるため、簡単化が生じる。これはまさに等価原理である。ゲージ理論の場合と同じく、それが可能となる領域は、上のラグランジアンの混合項を見るだけで見積もることができる。実際、式\eqref{eq:181}では、\(\pi\) と \(g_{\mu\nu}\) を混合する二次の項は \(\pi\) の運動項よりも微分の数が少なく、ある高いエネルギースケールを超えると無視できる。一般には、その答えはどの演算子が存在するかに依存する。ここでは、タドポール項だけが重要な最も簡単な場合、\(M_2=M_3=\cdots=0\) を考えよう。これは標準的なスローロール・インフレーションを含む。重力との主要な混合は、次の形の項から生じる。
\OriginalEquation{182}
次が成り立つことが分かる。
\OriginalEquation{183}
\trnote{式中の Kinetic term は運動項、Mixing term は混合項を表す。}

\(\delta g^{0i}\) は拘束変数であり、ある種の重力ポテンシャルに相当し、\(\pi\) によって決まる。短距離ではNewton近似が成り立つ。
\OriginalEquation{184}
したがって、
\OriginalEquation{185}
となる。最後から二番目の段階では空間微分について部分積分を行い、最後の段階では、エネルギーが \(E\) 程度のとき \(\partial_t\sim E\) であることを用いた。混合項は紫外領域で無視できる。これは一般相対論の等価原理である。混合を無視できるようになる実際のスケール \(E_{\mathrm{mix}}\) は、どの演算子を非零とするかに依存するが、エネルギー \(E\gg E_{\mathrm{mix}}\) では混合項を無視できることが保証される。

\(E\gg E_{\mathrm{mix}}\) の領域では、作用は劇的に簡単になり、
\OriginalEquation{186}
となる。
時間の微分同相変換が非線形に実現されることにより、\(\pi\) は非線形な「ブロック」
\OriginalEquation{187}
の形で現れなければならないことに注意しよう。これにより、異なる演算子の間に厳密な関係が与えられる。

あるインフレーション模型が与えられたとき、関心があるのは現在の宇宙論的観測に対する予言を計算することである。この観点から見ると、分離極限\eqref{eq:186}は、このような極端に赤外のスケールとはまったく関係がないように思われる。しかし、標準的な単一場のスローロール・インフレーションの場合と同様、通常の \(\zeta\) 変数という量が存在し、それが地平線の外側では摂動論の任意の次数で一定となることを証明できる。

したがって、問題は地平線通過の直後の相関関数を計算することに帰着する。このため、\(H\) 程度の赤外エネルギーカットオフを設けてラグランジアンを調べればよい。分離のスケール \(E_{\mathrm{mix}}\) が \(H\) より小さければ、\(\pi\) のラグランジアン\eqref{eq:186}は、\(E_{\mathrm{mix}}/H\) で抑制される項を除いて正しい予言を与える。そうでない場合にも、何か深刻なことが起こるわけではない。単に計量の揺らぎも保持すればよい。

これが、第2講で行った計算の正当化である。

すでに述べたように、ユニタリーゲージのラグランジアンの係数の時間変化はHubble時間に比べて緩やかである、すなわち何らかの一般化されたスローロール・パラメーターによって抑制されていると仮定している。このことは、係数のTaylor展開から生じる追加の \(\pi\) の項が小さいことを意味する。特に、三次の項のような関連演算子、すなわち赤外へ向かうにつれて支配的になる演算子は、スケール \(H\) では重要ではないため、ラグランジアン\eqref{eq:186}では無視した。それでも、文献\cite{Behbahani:2011it,Achucarro:2012sm}で行われているように、これらを取り入れることは容易である。

結論として、ラグランジアン\eqref{eq:186}を用いると、重力との混合が支配的ではないすべての観測量、例えば標準的なスローロール・インフレーションにおける非Gaussian性\cite{Maldacena:2002vr,Seery:2006vu}などを計算できる。ただし、スペクトルの傾きは、主要次数ではラグランジアン\eqref{eq:186}を用いて計算できる。先に見たように、その値は、混合項を無視して計算した地平線通過時のパワースペクトルから簡単に求められる。重力との混合が重要になり、Goldstoneボソンの作用だけでは予言に十分でなくなった場合でも、このアプローチの妥当性が失われるわけではないことを強調しておく。作用\eqref{eq:163}には模型に関するすべての情報が含まれており、重力との混合が大きい場合にも、あらゆる予言を計算するために用いることができる。

いくつかの点を強調しておこう。
\begin{itemize}
\item 上のラグランジアンは非常に簡単であり、単一の自由度をもつインフレーション模型をすべて統一的に記述する。
\item これは揺らぎの理論を記述する。そして、実際に検証しているのはその揺らぎである。
\item 素粒子物理学のカイラルラグランジアンに対応するものである。実際、実験の観点から見れば、インフレーションはGoldstoneボソンの理論であることを示している。
\item 単一の時計をもつインフレーション模型をすべて包含しているため、可能なシグナルについての定理を証明できる。
\item また、考えられるあらゆるシグナルを調べることができる。
\item 対称性によって何が強制され、どの演算子が許され、何が可能なのかが明確になる。例えば、\((\partial_i\pi)^2\) の係数は \(\dot H M_{\mathrm{Pl}}^2\) に固定される。\(\dot\pi^2\) についてはそうではない。このことは、微分の主要次数ではヌル・エネルギー条件を破れないことを意味する。\(\dot H>0\) であれば、\(\pi\) の空間的な運動項の符号が負のエネルギーに対応するものとなり、制御不能な不安定性をもたらす。EFTは、高階微分項を加えることによってこの問題をどのように解決できるかも教えてくれる。実際、現在知られているヌル・エネルギー条件（NEC）を破る方法は、すべてこの枠組みの中で見いだされたものである。
\item この形式は、標準模型を超える物理で通常行うことを実践するのに適している。対称性を加えて、ある演算子をほかの演算子に比べて増強することもできるし、特定の模型の紫外完備化を試みることもできる。
\item 明示的に揺らぎの理論であるため、演算子の重要性を容易に評価できる。例えば、スカラー場を用いた標準的な扱いでは、演算子 \((\partial\phi)^8\) は、\(\dot\phi_0^6(\partial\delta\phi)^2\) を通じて二次の作用に寄与する。この点はループ補正を調べるうえでも非常に有用である。揺らぎと微分の次数を固定すると、EFTには有限個の相殺項しか存在しないが、スカラー場の理論ではそうではない。実際、EFTの形式は、量子レベルでの \(\zeta\) の一定性を証明するために不可欠であった\cite{Pimentel:2012tw,Senatore:2012ya}。
\end{itemize}
